% !TeX program = pdflatex
% =============================================================================
% Lektionsplan (Diskussionsgrundlage, NICHT das Arbeitsblatt selbst)
% Lorentzkraft Teil 2: Freie Ladungen -- Kreisbahn und Anwendungen
% UZH Praktikum I, Sara Romer Urban, KS im Lee, HS2026
% Klassen: 4fh + 4cdeg
% Direkte Fortsetzung von Lorentzkraft Teil 1
% =============================================================================
\documentclass[11pt,a4paper]{article}

\usepackage[T1]{fontenc}
\usepackage[utf8]{inputenc}
\usepackage[ngerman]{babel}
\usepackage{lmodern}
\usepackage[a4paper,margin=2.2cm]{geometry}
\usepackage{array,booktabs,tabularx,xltabular}
\usepackage{enumitem}
\usepackage{xcolor}
\usepackage{amsmath}
\usepackage{siunitx}
\sisetup{per-mode=fraction,fraction-command=\frac}
\usepackage{url}
\Urlmuskip=0mu plus 1mu
\def\UrlBreaks{\do\/\do-\do_}

\definecolor{titleblue}{HTML}{183A6B}
\definecolor{accentblue}{HTML}{2E6F9E}
\definecolor{groupteal}{HTML}{1B7F72}
\definecolor{darkgray}{HTML}{4A4A4A}
\definecolor{warnred}{HTML}{9E2E2E}

\setlength{\parindent}{0pt}
\setlength{\parskip}{0.5em}

\newcommand{\Methode}[1]{{\small\color{groupteal}\textbf{#1}}}

\usepackage{titlesec}
\titleformat{\section}{\Large\bfseries\color{titleblue}}{\thesection}{0.6em}{}
\titlespacing*{\section}{0pt}{1.2em}{0.5em}

\begin{document}
\sloppy

\begin{center}
{\Huge\bfseries\color{titleblue}Lektionsplan: Lorentzkraft Teil 2}\\[0.2em]
{\Large\bfseries\color{titleblue}Freie Ladungen: Kreisbahn und Anwendungen}\\[0.4em]
{\large Doppellektion (2$\times$45 Min.) -- Diskussionsgrundlage}
\end{center}
\vspace{0.5em}

\noindent\textbf{Klassen:} 4fh + 4cdeg, gemeinsames Material -- Termin noch offen (spätere Doppellektion nach Teil 1)

\noindent\textbf{Lehrperson:} Loris Keller

\vspace{0.3em}
\noindent\textcolor{darkgray}{\textbf{Begriffe in diesem Dokument:} die
\textbf{Lorentzkraft} $F=qvB$ (Kraft auf eine einzelne bewegte Ladung,
hergeleitet in Teil 1) wirkt bei einer \emph{freien} Ladung im Magnetfeld
als \textbf{Zentripetalkraft} und erzeugt eine Kreisbahn mit Radius
$r=\frac{mv}{qB}$. Steht die Geschwindigkeit nicht senkrecht zum Feld,
entsteht stattdessen eine Schraubenbahn (Helix).}

\vspace{0.3em}
\noindent\textcolor{darkgray}{Dieses Dokument ist \textbf{keine}
Reinschrift eines Schülerdokuments, sondern die Planungsgrundlage dafür:
Lernziele, Aufbau der Lektion und die gewählte Didaktik-Methode pro
Phase -- analog zu `lorentzkraft-teil1-lektionsplan.tex`. Direkte
Fortsetzung von Lorentzkraft Teil 1: dort wurde die Kreisbahn im
Fadenstrahlrohr bereits \emph{qualitativ} erklärt (Predict-Observe-Explain,
Lorentzkraft als Zentripetalkraft); diese Frage wird hier \textbf{nicht}
wiederholt. Nach Absprache werden daraus lokal in VS Code ein
\textbf{Arbeitsblatt} (Theorie/Erarbeitung für diese Lektion), ein
\textbf{Aufgabenblatt} (zusätzliche Übungsaufgaben) und ein
\textbf{Gruppenpuzzle-Infomaterial} (sechs Kurztexte, je einer pro
Anwendung) erstellt -- das bestehende \textbf{Experimente-Dokument} aus
Teil 1 deckt das Fadenstrahlrohr bereits ab und wird für die
Radius-Variation in Teil 2 nur ergänzt, siehe Abschnitt 4.}

% =========================================================================
\section*{1. Abgrenzung und Lernziele}

\textbf{Im Fokus:} Teil 2 knüpft direkt an Teil 1 an, wo am
Fadenstrahlrohr bereits \emph{qualitativ} gezeigt und erklärt wurde, dass
eine freie, bewegte Ladung im Magnetfeld eine Kreisbahn beschreibt, weil
die Lorentzkraft $F=qvB$ stets senkrecht zur Geschwindigkeit steht und
damit als Zentripetalkraft wirkt (Teil 1, LZ6). \textbf{Diese Frage wird
hier nicht wiederholt.} Stattdessen baut Teil 2 direkt
\emph{quantitativ} darauf auf: aus $qvB=\frac{mv^2}{r}$ folgt der
Bahnradius $r=\frac{mv}{qB}$; kombiniert mit $v=\sqrt{2eU/m_e}$ (aus der
Beschleunigungsspannung) ergibt sich $r=\frac{1}{B}\sqrt{\frac{2m_eU}{e}}$
und daraus die historische Methode (Thomson) zur Bestimmung von $e/m_e$.
Danach folgt die Schraubenbahn als kurze Erweiterung (Geschwindigkeit
nicht senkrecht zum Feld), und abschliessend ein Schaufenster von sechs
Anwendungen desselben Grundprinzips, erarbeitet im Gruppenpuzzle.

\textbf{Bereits bekannt (nur Repetition, keine Neueinführung):} die
Lorentzkraft $F=qvB$ und ihre Richtungsbestimmung mit der
Drei-Finger-Regel (rechte Hand positiv, linke Hand negativ, Teil 1); der
Aufbau und die Funktionsweise des Fadenstrahlrohrs (Teil 1); dass die
Kreisbahn im Fadenstrahlrohr durch die Lorentzkraft als Zentripetalkraft
entsteht (Teil 1, LZ6, \emph{qualitativ} bereits gezeigt); die
Zentripetalkraft-Formel $F_Z=\frac{mv^2}{r}$ (Mechanik); die
Beschleunigungsspannung als Energieumwandlung $eU=\frac{1}{2}mv^2$
(Elektrizitätslehre, ergibt $v=\sqrt{2eU/m}$).

\textbf{Bewusst nicht Teil dieser Einheit (Kernstoff):} eine vollständige
mathematische Parametrisierung der Schraubenbahn (Helix-Gleichung in
Parameterform) -- es genügt die qualitative Zerlegung in $v_\parallel$
und $v_\perp$; relativistische Korrekturen bei hohen Geschwindigkeiten;
vertiefte technische Details einzelner Anwendungen über das im
Gruppenpuzzle vermittelte Mass hinaus (z.\,B. vollständige
Synchrotron-Physik) -- dort reicht das Grundprinzip plus ein
motivierender Ausblick.

\begin{itemize}[leftmargin=1.3cm,itemsep=0.3em]
  \item[LZ1] den Bahnradius $r=\frac{mv}{qB}$ herleiten, indem die
  Lorentzkraft $F=qvB$ (Teil 1) mit der bereits bekannten
  Zentripetalkraft $F_Z=\frac{mv^2}{r}$ gleichgesetzt wird -- als
  \textbf{quantitative} Fortsetzung der in Teil 1 bereits qualitativ
  gezeigten Kreisbahn (LZ6 dort), nicht als neue Einführung derselben.
  \item[LZ2] mithilfe der Beschleunigungsspannung
  ($v=\sqrt{2eU/m_e}$) die Radiusformel zu
  $r=\frac{1}{B}\sqrt{\frac{2m_eU}{e}}$ umformen und daraus
  $e/m_e=\frac{2U}{r^2B^2}$ bestimmen -- die historische Methode
  (Thomson) zur Messung des Verhältnisses $e/m_e$.
  \item[LZ3] anhand von $qvB=\frac{mv^2}{r}$ begründet vorhersagen, wie
  sich der Bahnradius bei veränderter Beschleunigungsspannung bzw.
  Feldstärke verhält, die Vorhersage am Fadenstrahlrohr überprüfen und
  das Ergebnis mit LZ1/LZ2 erklären (Predict-Observe-Explain).
  \item[LZ4] die Schraubenbahn (Helix) erklären, die entsteht, wenn die
  Geschwindigkeit nicht senkrecht zum Feld steht: $v_\parallel$ bleibt
  unbeeinflusst (gleichförmige Bewegung), $v_\perp$ erzeugt die
  Kreisbewegung -- zusammen ergibt das eine Schraubenbahn.
  \item[LZ5] mindestens eine der sechs Anwendungen (Massenspektrograph,
  Geschwindigkeitsfilter/Wien-Filter, Zyklotron, Hall-Effekt,
  MHD-Antrieb, Erdmagnetfeld) im Detail erklären können und im
  Gruppenpuzzle einer neu gemischten Gruppe verständlich vermitteln.
  \item[LZ6] erkennen und begründen, dass alle sechs Anwendungen auf
  demselben Grundprinzip beruhen -- Lorentzkraft als Zentripetalkraft
  (Massenspektrograph, Zyklotron, Erdmagnetfeld) bzw.\ als
  Kräftegleichgewicht bei gekreuzten Feldern (Wien-Filter, Hall-Effekt) --,
  angewandt auf ein je anderes Gerät oder Phänomen.
\end{itemize}

% =========================================================================
\section*{2. Aufbau der Doppellektion}

\begin{xltabular}{\textwidth}{@{}p{2.6cm}X p{1.2cm}@{}}
\toprule
\textbf{Phase} & \textbf{Inhalt \& Methode} & \textbf{Zeit} \\
\midrule
\endhead

Einstieg &
\textbf{Rückblick + neue quantitative Frage.} Kurzer Rückblick auf Teil 1
(\glqq Wir wissen schon: Elektronen im Fadenstrahlrohr laufen auf einer
Kreisbahn, weil die Lorentzkraft als Zentripetalkraft wirkt.\grqq{}),
danach die neue Leitfrage: Wovon hängt der Radius dieser Kreisbahn ab?
Lehrervortrag mit kurzem Unterrichtsgespräch, Plenum. &
4' \\[0.4em]

POE &
\textbf{Predict-Observe-Explain: Bahnradius am Fadenstrahlrohr
(quantitativ).} \emph{Predict:} die SuS sagen anhand von
$qvB=\frac{mv^2}{r}$ voraus, wie sich $r$ bei erhöhter Spannung $U$ bzw.\
erhöhter Feldstärke $B$ verändert. \emph{Observe:} $U$ bzw.\ $B$ am
Fadenstrahlrohr tatsächlich variieren und die Radiusänderung
beobachten. \emph{Explain:} Herleitung von $r=\frac{mv}{qB}$ aus der
Zentripetalkraft-Gleichung, damit die beobachtete Abhängigkeit erklären
(LZ1, LZ3). \Methode{M4 Predict-Observe-Explain} -- Lehrer-Demonstration,
Plenum. &
12' \\[0.4em]

Vertiefung &
\textbf{Historische $e/m_e$-Bestimmung.} Direkte Fortsetzung der
Herleitung: mit $v=\sqrt{2eU/m_e}$ (aus der Beschleunigungsspannung) wird
$r=\frac{1}{B}\sqrt{2m_eU/e}$ zu $e/m_e=\frac{2U}{r^2B^2}$ umgeformt --
die historische Methode (Thomson) zur Bestimmung von $e/m_e$ (LZ2).
Gemeinsam entwickelte Umformung im Plenum. &
8' \\[0.4em]

Übung &
\textbf{Anwendung von $r=mv/qB$ und $e/m_e$.} Kurze Rechenaufgaben mit
konkreten Zahlenwerten (z.\,B.\ aus einer typischen
Fadenstrahlrohr-Apparatur). Individuelles Üben. &
8' \\[0.4em]

Theorie &
\textbf{Schraubenbahn.} Was passiert, wenn $v$ schräg zu $B$ steht?
Zerlegung in $v_\parallel$ (unbeeinflusst) und $v_\perp$ (erzeugt die
Kreisbewegung) -- zusammen ergibt das eine Schraubenbahn (LZ4).
Lehrervortrag mit Skizze, Plenum. &
5' \\[0.4em]

Einstieg &
\textbf{Advance Organizer: Bandbreite der Anwendungen.} Kurzer Überblick,
der zeigt, wie breit gestreut die sechs Anwendungen sind -- z.\,B.\ ein
Bild der Yamato-1 (MHD-Antrieb) als neuer Aufhänger, plus bewusster
Rückverweis auf das Polarlicht-Beispiel aus dem Advance Organizer von
Teil 1: \glqq Erinnert ihr euch an das Polarlicht-Beispiel von letzter
Doppellektion? Genau das ist eine der sechs Anwendungen, die wir jetzt
vertiefen.\grqq{} Details siehe Abschnitt 2.1. Lehrervortrag, kurz,
Plenum. &
3' \\[0.4em]

Gruppenpuzzle &
\textbf{Erarbeitung (Expertengruppen).} Je eine Gruppe erarbeitet
anhand vorbereiteten Info-Materials eine der sechs Anwendungen
(Massenspektrograph, Geschwindigkeitsfilter/Wien-Filter, Zyklotron,
Hall-Effekt, MHD-Antrieb, Erdmagnetfeld) -- alle bauen auf demselben
Grundprinzip auf (LZ5, LZ6). \Methode{Partner-/Gruppenpuzzle} --
Gruppenarbeit. &
14' \\[0.4em]

Gruppenpuzzle &
\textbf{Vermittlung (neu gemischte Gruppen).} Neu gemischte Gruppen (je
ein Mitglied pro Anwendung), jede Person erklärt ihre Anwendung den
anderen (LZ5). \Methode{Partner-/Gruppenpuzzle} -- Gruppenarbeit. &
14' \\[0.4em]

Sicherung &
\textbf{Plenum: gemeinsames Grundprinzip.} Kurzer Kurzcheck/Quiz zu allen
sechs Anwendungen, mit Fokus darauf, dass alle auf demselben
Grundprinzip beruhen -- Lorentzkraft als Zentripetalkraft bzw.\ als
Kräftegleichgewicht bei gekreuzten Feldern (LZ6). Plenum. &
7' \\[0.4em]

Abschluss &
\textbf{Zusammenfassung der gesamten Lorentzkraft-Einheit.} Kurzer
Rückblick von der Kraft auf den stromdurchflossenen Leiter (Teil 1) über
die einzelne freie Ladung bis zu den sechs Anwendungen (Teil 2). Plenum. &
3' \\

\bottomrule
\end{xltabular}

\textcolor{darkgray}{\small Summe: 78 Min. -- 12 Min. Puffer innerhalb der
90 Min. der Doppellektion -- siehe Zeitbudget in Abschnitt 5, das
Gruppenpuzzle (28' zusammen) ist der grösste Zeitblock und läuft am
ehesten Gefahr, sich zu verlängern.}

\subsection*{2.1 Ausgestaltungsvorschlag: Advance Organizer vor dem Gruppenpuzzle}

Ziel: in 3 Minuten die Bandbreite der sechs Anwendungen zeigen und Lust
auf das Gruppenpuzzle machen, ohne die Inhalte selbst vorwegzunehmen.

\textbf{Bild 1 (neu): Yamato-1.} Das japanische Versuchsschiff Yamato-1
fährt ohne Propeller -- ein magnetohydrodynamischer Antrieb beschleunigt
das Meerwasser direkt über die Lorentzkraft. Frage an die Klasse: \glqq
Wie kann ein Schiff fahren, ohne dass sich irgendetwas dreht?\grqq{}

\textbf{Bild 2 (Rückverweis auf Teil 1): Polarlicht.} \glqq Erinnert ihr
euch an das Polarlicht-Beispiel vom Advance Organizer letzte
Doppellektion? Genau das ist eine der sechs Anwendungen, die wir jetzt im
Detail vertiefen.\grqq{} Der Rückverweis zeigt, dass der damalige
Aufhänger jetzt mit vollem Handwerkszeug erklärbar wird.

\textbf{Überleitung:} \glqq Vom Massenspektrographen über das Zyklotron
bis zu den Polarlichtern -- ganz unterschiedliche Geräte und Phänomene,
aber alle beruhen auf demselben Prinzip: der Lorentzkraft. In den
nächsten Gruppen werdet ihr je eine dieser sechs Anwendungen zu
Expert:innen und erklärt sie danach den anderen.\grqq{}

\subsection*{2.2 Kerninhalte und Leitfragen je Phase}

\begin{itemize}[leftmargin=1.3cm,itemsep=0.3em]
  \item \textbf{Einstieg (Rückblick + neue Frage):} die Kreisbahn ist aus
  Teil 1 bekannt -- jetzt geht es um den Radius. \emph{Leitfrage:} Wovon
  hängt der Radius der Kreisbahn ab?
  \item \textbf{POE Fadenstrahlrohr (quantitativ):} $r=mv/qB$, hergeleitet
  aus der Zentripetalkraft. \emph{Leitfrage:} Wie verändert sich der
  Radius, wenn ich Spannung oder Feldstärke ändere?
  \item \textbf{Vertiefung $e/m_e$-Bestimmung:} mit $v=\sqrt{2eU/m_e}$
  wird $e/m_e$ aus $U$, $r$, $B$ bestimmbar. \emph{Leitfrage:} Wie lässt
  sich das Verhältnis $e/m_e$ aus der Kreisbahn bestimmen?
  \item \textbf{Übung:} Anwendung von $r=mv/qB$ und der $e/m_e$-Formel an
  Beispielen. \emph{Leitfrage:} Wie rechne ich mit diesen Formeln sicher?
  \item \textbf{Theorie Schraubenbahn:} $v_\parallel$ bleibt
  unbeeinflusst, $v_\perp$ erzeugt die Kreisbewegung -- zusammen eine
  Helix. \emph{Leitfrage:} Was passiert, wenn die Geschwindigkeit schräg
  zum Feld steht?
  \item \textbf{Advance Organizer Gruppenpuzzle:} dasselbe Grundprinzip
  taucht in ganz unterschiedlichen Geräten/Phänomenen auf.
  \emph{Leitfrage:} Was haben ein Massenspektrograph, ein Zyklotron und
  die Polarlichter gemeinsam?
  \item \textbf{Gruppenpuzzle Erarbeitung:} je eine der sechs Anwendungen
  im Detail verstehen. \emph{Leitfrage:} Wie funktioniert \emph{meine}
  Anwendung genau?
  \item \textbf{Gruppenpuzzle Vermittlung:} die eigene Anwendung
  verständlich weitergeben. \emph{Leitfrage:} Kann ich meine Anwendung so
  erklären, dass es die anderen verstehen?
  \item \textbf{Sicherung/Plenum:} alle sechs Anwendungen beruhen auf
  demselben Grundprinzip. \emph{Leitfrage:} Was ist allen sechs
  Anwendungen gemeinsam?
  \item \textbf{Abschluss:} Zusammenfassung der ganzen
  Lorentzkraft-Einheit (Teil 1 + Teil 2). \emph{Leitfrage:} Was haben wir
  in den letzten beiden Doppellektionen über die Lorentzkraft gelernt?
\end{itemize}

% =========================================================================
\section*{3. Begründung der Methodenwahl}

Der Einstieg verzichtet bewusst auf einen erneuten
Predict-Observe-Explain-Zyklus zur \emph{qualitativen} Kreisbahn-Frage --
die ist bereits in Teil 1 beantwortet. Stattdessen wird die
Predict-Observe-Explain-Methode hier auf eine \emph{neue}, quantitative
Frage angewendet (Abhängigkeit des Radius von $U$ und $B$), bei der die
SuS dank der in Teil 1 bereits gelegten Grundlage (Lorentzkraft als
Zentripetalkraft, Zentripetalkraft-Formel aus der Mechanik) tatsächlich
eine begründete Vorhersage treffen können -- ohne diese Grundlage wäre
auch diese Vorhersage nur Raten.

Die $e/m_e$-Bestimmung folgt bewusst \emph{direkt} auf die
Radius-Herleitung, nicht als separater Exkurs später: sie ist nur eine
Umformung derselben Formel um eine zusätzliche Beziehung
($v=\sqrt{2eU/m_e}$), sodass die SuS sehen, wie unmittelbar diese
historische Messmethode aus der bereits gemachten Herleitung folgt.

Für die sechs Anwendungen eignet sich das Gruppenpuzzle (Jigsaw)
besonders gut, weil die sechs Themen inhaltlich unabhängig voneinander
und ungefähr gleich aufwendig sind -- jede Gruppe kann sich auf eine
Anwendung konzentrieren, statt dass die Lehrperson alle sechs im
Frontalunterricht vorstellen muss. Der Advance Organizer davor motiviert
die Gruppenarbeit, indem er zeigt, wie unterschiedlich (und
alltagsrelevant) die sechs Anwendungen sind, und schafft mit dem
Polarlicht-Rückverweis eine bewusste Klammer zu Teil 1.

% =========================================================================
\section*{4. Anschlussdokumente: Arbeitsblatt, Aufgabenblatt, Gruppenpuzzle-Material}

Aus diesem Lektionsplan entstehen drei weitere Dokumente (das
Experimente-Dokument aus Teil 1 deckt das Fadenstrahlrohr bereits ab und
wird nur um die Radius-Variation ergänzt):

\begin{itemize}[leftmargin=1.2cm,itemsep=0.4em]
  \item \textbf{Arbeitsblatt:} Theorie- und Erarbeitungsdokument für
  diese Lektion (analog zu Teil 1), entlang der in Abschnitt 2
  beschriebenen Phasen.
  \item \textbf{Aufgabenblatt:} zusätzliche Übungsaufgaben zu $r=mv/qB$,
  zur $e/m_e$-Bestimmung und zur Schraubenbahn.
  \item \textbf{Gruppenpuzzle-Infomaterial:} sechs vergleichbar
  umfangreiche Kurztexte/Arbeitsblätter, je einer pro Anwendung
  (Massenspektrograph, Geschwindigkeitsfilter/Wien-Filter, Zyklotron,
  Hall-Effekt, MHD-Antrieb, Erdmagnetfeld) -- muss vor der Durchführung
  noch erstellt werden.
\end{itemize}

% =========================================================================
\section*{5. Offene Punkte zur Besprechung}

\begin{itemize}[leftmargin=1.2cm,itemsep=0.4em]
  \item \textbf{Zeitbudget:} Zehn Phasen (Einstieg, POE, Vertiefung
  $e/m_e$, Übung, Schraubenbahn, Advance Organizer, Gruppenpuzzle-
  Erarbeitung, Gruppenpuzzle-Vermittlung, Sicherung, Abschluss) sind bei
  12' Puffer ambitioniert für eine Doppellektion -- insbesondere das
  Gruppenpuzzle (28' zusammen) läuft am ehesten Gefahr, sich zu
  verlängern. Erster Streichkandidat bei Zeitdruck: die Übungsphase
  kürzen oder als Hausaufgabe auslagern.
  \item \textbf{Info-Material für das Gruppenpuzzle:} muss für alle
  sechs Anwendungen noch erstellt werden (Kurztexte/Arbeitsblätter pro
  Anwendung) -- Umfang und Schwierigkeitsgrad sollten über alle sechs
  vergleichbar sein, damit keine Expertengruppe deutlich mehr/weniger
  Zeit braucht.
  \item \textbf{Auswahl/Reihenfolge der sechs Anwendungen:} passt die
  Auswahl (Massenspektrograph, Wien-Filter, Zyklotron, Hall-Effekt,
  MHD-Antrieb, Erdmagnetfeld) so, oder gibt es eine, die eher wegfallen
  oder ersetzt werden sollte (z.\,B.\ falls Hall-Effekt bereits an
  anderer Stelle im Lehrplan vorkommt)?
  \item \textbf{Schraubenbahn -- Tiefe:} reicht die qualitative
  Behandlung (Zerlegung in $v_\parallel$/$v_\perp$), oder soll es eine
  eigene kleine Rechenaufgabe dazu geben?
  \item \textbf{Technische Voraussetzung:} die Predict-Observe-Explain-
  Phase braucht ein Fadenstrahlrohr mit regelbarer
  Beschleunigungsspannung \emph{und} regelbarem Feld (Helmholtzspulen mit
  einstellbarem Strom) -- bitte gegenprüfen, ob das vorhandene Gerät
  beides erlaubt.
  \item \textbf{Sozialform Gruppenpuzzle:} Gruppengrösse und genaue
  Durchführung (z.\,B.\ sechs Fünfergruppen zu Beginn, danach sechs neu
  gemischte Fünfergruppen mit je einem Mitglied pro Anwendung) noch
  festzulegen, abhängig von der Klassengrösse.
  \item \textbf{Termin:} Datum der Teil-2-Doppellektion noch offen --
  abhängig von der Terminierung von Teil 1 (15./18.09.2026) und dem
  weiteren Stoffplan.
\end{itemize}

\end{document}
